Table~\ref{tab:decay_ablation} shows that the constant kernel already
improves the shared TS-only baseline. Its mean QLIKE is 0.59249,
a reduction of 7.16\%, and its mean log-MSE is 0.80461,
a reduction of 0.63\%. A model without within-window discounting
can therefore improve forecast accuracy. The evidence does not
support treating temporal decay as a necessary condition for any
benefit from the news correction.

All three decaying specifications improve further on both average
losses. The refitted exponential model reduces QLIKE by 10.57\%
and log-MSE by 2.48\% relative to TS-only. Power-law decay yields
a slightly larger QLIKE reduction, 10.94\%, and a slightly smaller
log-MSE reduction, 2.31\%. Linear decay has the lowest losses in
this comparison: QLIKE is 0.54907 and log-MSE is 0.78077,
corresponding to reductions of 13.96\% and 3.58\%. It also has
the highest mean $R^2$, 0.440, and DA, 0.725.

The comparison supports a role for temporal discounting beyond a
constant aggregation of article magnitudes, while qualifying a
stronger claim that the gain disappears without decay. It also
shows that functional form can matter quantitatively. Exponential
decay offers a smooth response and a direct half-life interpretation,
but it is not the most accurate kernel in this experiment. Since
each specification is estimated with one fixed seed, these rankings
describe the fitted models rather than establishing an ordering
that is invariant to optimization variability.
